% =====================================================================
%  Tiền xử lý: Tách câu thành Unit Opinion Sentence bằng LLM
%  — \input{} vào ngay trước \subsection{Kết quả đánh giá với tách câu}
% =====================================================================

\subsection{Phương pháp tách câu bằng mô hình ngôn ngữ lớn (UOS)}
\label{subsec:uos_method}

\subsubsection*{Động lực}

Một câu đánh giá khách sạn thường đề cập đến nhiều khía cạnh cùng
lúc trong một mệnh đề duy nhất, ví dụ:

\begin{quote}
\textit{``Overall, the quality of the restaurant, service, rooms as well as
the staff all bring a great feeling, well worth the experience.''}
\end{quote}

Câu trên chứa bốn khía cạnh độc lập: \textit{restaurant},
\textit{service}, \textit{rooms}, \textit{staff}. Khi đưa nguyên câu
vào mô hình ABSA, biểu diễn ngữ nghĩa của một khía cạnh bị pha trộn
bởi các khía cạnh còn lại, dẫn đến hiện tượng bỏ sót dự đoán
(\textit{under-prediction}).

Để khắc phục, bước tiền xử lý tách mỗi câu gốc thành các
\textbf{Unit Opinion Sentence (UOS)} — những câu hoàn chỉnh, tự thân,
mỗi câu chỉ chứa đúng một ý kiến về đúng một khía cạnh.

\subsubsection*{Kiến trúc pipeline}

Pipeline tách UOS giao tiếp với mô hình ngôn ngữ thông qua
\textbf{Ollama API} (endpoint \texttt{/api/chat}) chạy cục bộ với mô
hình \textbf{Qwen3-8B}. Luồng xử lý từng câu gồm bốn bước:

\begin{enumerate}
  \item \textbf{Xây dựng prompt} (\texttt{prompt.py}) — Nhúng câu đầu
        vào vào prompt cố định \textit{linguistic\_v11} chứa định nghĩa
        aspect term, quy tắc tách và 14 ví dụ few-shot.
  \item \textbf{Gọi LLM} — Gửi prompt tới Qwen3-8B qua Ollama;
        tham số \texttt{think=False} tắt chuỗi suy luận nội bộ để giảm
        độ trễ.
  \item \textbf{Phân tích đầu ra} (\texttt{parsing.py}) — Trích xuất
        mảng JSON từ văn bản thô, loại bỏ code-fence (\texttt{```}) và
        thẻ \texttt{<think>}, chuẩn hóa dấu câu cuối mỗi UOS.
  \item \textbf{Kiểm tra hợp lệ} (\texttt{validation.py}) — Phát hiện
        tách quá mức (\textit{over-split}), cặp tương phản chưa tách,
        hoặc token không xuất hiện trong câu gốc.
\end{enumerate}

\subsubsection*{Thiết kế prompt}

Prompt \textit{linguistic\_v11} được thiết kế theo hướng quy tắc ngôn
ngữ, không phân loại aspect hay sentiment trực tiếp.
Bảng~\ref{tab:uos_rules} tóm tắt các quy tắc tách/không tách chính.

\begin{table}[H]
\centering
\caption{Quy tắc tách câu trong prompt \textit{linguistic\_v11}}
\label{tab:uos_rules}
\small
\begin{tabular}{p{0.46\linewidth} p{0.46\linewidth}}
\toprule
\textbf{TÁCH} & \textbf{KHÔNG TÁCH} \\
\midrule
Nhiều aspect term, mỗi cái mang ý kiến riêng &
Nhiều tính từ bổ nghĩa cho cùng một aspect term \\[4pt]
Danh sách aspect nối bằng dấu phẩy hoặc ``and'' (có thể tách độc lập) &
Cặp danh từ ghép chặt thành một khái niệm
(vd.\ \textit{staff and customer service}) \\[4pt]
Tính từ dùng chung phân phối lại cho từng aspect &
Câu logistic/bối cảnh kèm một ý kiến duy nhất \\[4pt]
Hai ý kiến tương phản trên một aspect &
Phát biểu sự kiện không mang cảm xúc \\
\bottomrule
\end{tabular}
\end{table}

Ngoài ra mỗi UOS phải tuân thủ: (i)~chỉ dùng từ ngữ có trong câu gốc
— không diễn giải, không thay thế; (ii)~thêm ``The''/``A'' làm chủ
ngữ nếu thiếu; (iii)~thêm copula (``is'', ``was'') nếu câu bị phân
mảnh.

\subsubsection*{Thống kê tách câu trên tập kiểm tra}

Bảng~\ref{tab:uos_stats} và Bảng~\ref{tab:uos_dist} tổng hợp kết quả
tách UOS trên toàn bộ 312 câu của tập kiểm tra.

\begin{table}[H]
\centering
\caption{Tổng quan tách UOS trên tập kiểm tra}
\label{tab:uos_stats}
\begin{tabular}{@{}lr@{}}
\toprule
\textbf{Chỉ số} & \textbf{Giá trị} \\
\midrule
Tổng câu đầu vào               & 312 \\
Tổng UOS đầu ra                & 612 \\
Số UOS trung bình / câu        & 1{,}96 \\
Câu được tách (số UOS $> 1$)  & 174 (55{,}8\%) \\
Câu giữ nguyên (số UOS $= 1$) & 138 (44{,}2\%) \\
\bottomrule
\end{tabular}
\end{table}

\begin{table}[H]
\centering
\caption{Phân phối số UOS đầu ra trên 312 câu đầu vào}
\label{tab:uos_dist}
\begin{tabular}{@{}crr@{}}
\toprule
\textbf{Số UOS} & \textbf{Số câu} & \textbf{Tỉ lệ (\%)} \\
\midrule
1 & 138 & 44{,}2 \\
2 &  95 & 30{,}4 \\
3 &  47 & 15{,}1 \\
4 &  22 &  7{,}1 \\
5 &   6 &  1{,}9 \\
6 &   3 &  1{,}0 \\
7 &   1 &  0{,}3 \\
\midrule
\textbf{Tổng} & \textbf{312} & \textbf{100{,}0} \\
\bottomrule
\end{tabular}
\end{table}

Kết quả cho thấy gần 56\% câu đầu vào chứa từ hai khía cạnh trở lên
và được tách thành nhiều UOS. Phổ biến nhất là tách thành 2 UOS
(30,4\%), phản ánh đặc thù dữ liệu đánh giá khách sạn: người dùng
thường đề cập đồng thời nhiều khía cạnh trong một câu ngắn.

Bảng~\ref{tab:uos_examples} minh họa các trường hợp tách tiêu biểu.

\begin{table}[H]
\centering
\caption{Ví dụ tách câu gốc thành UOS (trích từ tập kiểm tra)}
\label{tab:uos_examples}
\small
\begin{tabular}{p{0.07\linewidth} p{0.40\linewidth} p{0.43\linewidth}}
\toprule
\textbf{} & \textbf{Câu gốc} & \textbf{Các UOS đầu ra} \\
\midrule
$1{\to}1$ &
\textit{they upgraded us to a beautiful suite without us even asking} &
(1)~They upgraded us to a beautiful suite without us even asking. \\[6pt]

$1{\to}2$ &
\textit{The service staff is very attentive and conscientious} &
(1)~The service staff is very attentive. \newline
(2)~The service staff is very conscientious. \\[6pt]

$1{\to}3$ &
\textit{xe takes you to your door, enthusiastic driver, cute cheese sticks} &
(1)~Xe takes you to your door. \newline
(2)~The driver is enthusiastic. \newline
(3)~The cheese sticks are cute. \\[6pt]

$1{\to}4$ &
\textit{Overall, the quality of the restaurant, service, rooms as well as the staff all bring a great feeling, well worth the experience.} &
(1)~The restaurant is great. \newline
(2)~The service is great. \newline
(3)~The rooms are great. \newline
(4)~The staff is great. \\
\bottomrule
\end{tabular}
\end{table}

Trường hợp $1{\to}4$ minh họa khả năng mô hình xác định các aspect
term ẩn, phân phối lại vị từ chung ``is great'' cho từng UOS, đồng
thời bỏ qua phần bình luận tổng quan (``overall'', ``well worth the
experience'') vì không gắn với một khía cạnh cụ thể.
Sau bước tách UOS, mỗi UOS là một mẫu độc lập trong pipeline ABSA,
đảm bảo mỗi câu đầu vào chỉ chứa đúng một cặp (aspect term, sentiment)
và giảm hiện tượng bỏ sót dự đoán ở giai đoạn phân loại.
